\documentclass[10pt,a4paper]{amsart}
\usepackage[margin=19mm]{geometry}
\usepackage{amsmath,amssymb,mathtools}
\usepackage[scr=rsfs,cal=euler]{mathalfa}
\usepackage{xfrac}
\usepackage[unicode,hidelinks,bookmarks=false]{hyperref}
\newcommand{\ie}{i.e.\ }
\newcommand{\cf}{cf.\ }
\newcommand{\resp}{respectively\ }
\newcommand{\Subsection}{\subsection}
\providecommand{\nobreakdash}{\nobreak\hbox{-}}
\setlength{\emergencystretch}{2em}
\multlinegap0pt
\usepackage{bm}
\usepackage{bbm}
\usepackage{dsfont}
\usepackage{xspace}
\usepackage{cancel}
\usepackage{mathtools}
\usepackage{amsthm}
\makeatletter
\@ifundefined{Theorem}{\newtheorem{Theorem}{Theorem}}{}
\@ifundefined{Lemma}{\newtheorem{Lemma}[Theorem]{Lemma}}{}
\makeatother
\DeclareMathOperator{\sgn}{sgn}
\renewcommand{\geq}{\geqslant}
\renewcommand{\leq}{\leqslant}
\title[Affine Jacobi-Trudi Identities and $q,t$-Rogers-Ramanujan Id]{Affine Jacobi-Trudi Identities and $q,t$-Rogers-Ramanujan Identities}
\author{S. Ole Warnaar}
\date{Source: arXiv:2511.17034v2 (CC BY 4.0)}
\begin{document}
\maketitle

\noindent\emph{Source and licence.} Affine Jacobi-Trudi Identities and $q,t$-Rogers-Ramanujan Identities. Version arXiv:2511.17034v2 (CC BY 4.0), distributed under the Creative Commons Attribution 4.0 International licence, as recorded in the arXiv licence metadata for this exact version and shown on the abstract page https://arxiv.org/abs/2511.17034v2. Full rights notice: licenses/arxiv-2511.17034-CC-BY-4.0.txt.\par\smallskip

\noindent\textit{Editorial scope.} Counted unit: the Lemma whose source label is \texttt{Lem\_determinant} in the article (source0.tex lines 1487--1518, source label \texttt{Lem\_determinant}), delivered together with the article's own complete original proof of it (source0.tex lines 1521--1666, one \texttt{proof} environment of 146 lines). No mathematics is written, repaired or re-derived: statement and proof are the article's own text line for line. Required context retained uncounted: Eq\_ee (lines 1483--1485). Every definition, display and label of those lines is retained unchanged. Omitted from this package and not redistributed: 1--1482, 1486--1486, 1519--1520, 1667--4641. Nothing inside a retained excerpt is omitted or altered: every retained body is delivered exactly as the article's lines are written, so no omission receipt of any kind accompanies this package. Citations: the retained excerpts carry no citation command at all, so no bibliography entry of the article is transcribed and no reference is redistributed. Reference adaptation: none was needed and none was made; every reference inside the delivered unit resolves inside it, so no unresolved reference or citation is produced. Printed slips kept literal: no printed word, symbol, hypothesis or number of the counted statement or of its proof was altered. Layout and class: the article's own class is \texttt{sigma} with packages including ifpdf, cases, mathdots, fouridx, tikz; the wrapper is the lane's pinned \texttt{amsart} editorial wrapper at 10pt on A4, which restates the theorem-like environments the retained text needs and adds the article's own macro definitions that the retained text needs (\texttt{\textbackslash geq}, \texttt{\textbackslash leq}, \texttt{\textbackslash sgn}), with extra wrapper packages (bm, bbm, dsfont, xspace, cancel, mathtools). The delivered numbering is the unit's own. Assets: no retained span contains a figure environment, a graphics-inclusion command, a PDF-page inclusion, a table or a TikZ picture; no image file is copied into this package, the package declares no asset dependency and no image file is redistributed with it. Licence: version arXiv:2511.17034v2 of this article is distributed under the Creative Commons Attribution 4.0 International licence, as recorded in the arXiv licence metadata for this exact version and shown on the abstract page https://arxiv.org/abs/2511.17034v2. A full rights notice is delivered at licenses/arxiv-2511.17034-CC-BY-4.0.txt. Characters: every retained excerpt is pure ASCII, so the delivered engine, XeLaTeX with its default Latin Modern text font, reports no missing character.
\begin{equation}\label{Eq_ee}
\ddot{e}_r(x)=\dot{e}_r(x)+\dot{e}_{r-1}(x).
\end{equation}

\begin{Lemma}\label{Lem_determinant}
For $K$ an arbitrary integer,
\begin{subequations}
\begin{align}
\begin{aligned}[b]
&\sum_{y\in\mathbb{Z}^k}
\det_{1\leq i,j\leq k}\bigl(u^{y_i}
t^{\frac{1}{2}Ky_i^2-(j-\frac{1}{2})y_i}
 (\dot{e}_{n-i+j-Ky_i}(x)+\dot{e}_{n+i+j-Ky_i-1}(x) )\bigr)
  \\
&\qquad{}=\frac{1}{2}\sum_{y\in\mathbb{Z}^k}
\det_{1\leq i,j\leq k}\bigl(u^{y_i}
t^{\frac{1}{2}Ky_i^2-(j-\frac{1}{2})y_i}
 (\ddot{e}_{n-i+j-Ky_i+1}(x)+\ddot{e}_{n+i+j-Ky_i-1}(x) )\bigr)
 \end{aligned}
\label{Eq_doteddote}
\end{align}
and
\begin{align}
&\sum_{y\in\mathbb{Z}^k}
\det_{1\leq i,j\leq k}\bigl(u^{y_i}
t^{\frac{1}{2}Ky_i^2-j y_i} (
\dot{e}_{n-i+j-Ky_i}(x)-\dot{e}_{n+i+j-Ky_i}(x) )\bigr) \notag \\
&\qquad{}=
\sum_{y\in\mathbb{Z}^k}
\det_{1\leq i,j\leq k}\bigl(u^{y_i}
t^{\frac{1}{2}Ky_i^2-j y_i} (
\ddot{e}_{n-i+j-Ky_i+1}(x)-\ddot{e}_{n+i+j-Ky_i}(x) )\bigr).
\label{Eq_doteddote2}
\end{align}
\end{subequations}
\end{Lemma}

\begin{proof}
First we consider \eqref{Eq_doteddote}.
Fix $y\in\mathbb{Z}^k$ such that $y_1\geq y_2\geq\cdots\geq y_k$.
Then \eqref{Eq_doteddote} is a~consequence~of
\begin{align}
&\sum_{S_k\cdot y} \det_{1\leq i,j\leq k}\bigl( t^{-jy_i}
 (\dot{e}_{n-i+j-Ky_i}(x)+\dot{e}_{n+i+j-Ky_i-1}(x) )\bigr) \notag \\
&\qquad=\frac{1}{2}\sum_{S_k\cdot y} \det_{1\leq i,j\leq k}\bigl(
t^{-jy_i}
 (\ddot{e}_{n-i+j-Ky_i+1}(x)+\ddot{e}_{n+i+j-Ky_i-1}(x) )\bigr),
\label{Eq_stronger}
\end{align}
where $S_k\cdot y$ denotes the $S_k$ orbit of $y$.
Summing over all $k!$ permutations of $y$ instead of $S_k\cdot y$
amounts to multiplying \eqref{Eq_stronger} by the size of
the stabilizer of $y$, and in the following it will
be more convenient to consider \eqref{Eq_stronger} with $y$ replaced by
$w(y)=(y_{w_1},\dots,y_{w_k})$ and the sum over $S_k\cdot y$ replaced
by a sum over $w\in S_k$.
Then, after applying \eqref{Eq_ee} to the right-hand side, the identity
to be proved becomes the special case --
replace $(x_i,z_i)\mapsto (t^{-y_i},Ky_i)$ followed
by~${E_r\mapsto \dot{e}_{n-r}(x)}$~-- of the formal identity
\begin{align*}
&\sum_{w\in S_k} \det_{1\leq i,j\leq k}\bigl(x_{w_i}^j
\bigl(E_{z_{w_i}+i-j}+E_{z_{w_i}-i-j+1}\bigr)\bigr) \\
&\qquad{}=\frac{1}{2}\sum_{w\in S_k}
\det_{1\leq i,j\leq k}\bigl(
x_{w_i}^j \bigl(E_{z_{w_i}+i-j-1}+E_{z_{w_i}+i-j}+
E_{z_{w_i}-i-j+1}+E_{z_{w_i}-i-j+2}\bigr)\bigr).
\end{align*}
Dispensing with the determinants, this is the same as
\[
\sum_{\sigma,w\in S_k}\sgn(\sigma)\prod_{i=1}^k
x_{w_i}^{\sigma_i} F_{i,i}(\sigma,w)
=\frac{1}{2}\sum_{\sigma,w\in S_k}\sgn(\sigma)\prod_{i=1}^k
x_{w_i}^{\sigma_i} (F_{i,i-1}(\sigma,w)+F_{i,i}(\sigma,w) ),
\]
where
\[
F_{i,j}(\sigma,w):=E_{z_{w_i}+j-\sigma_i}+E_{z_{w_i}-j-\sigma_i+1}.
\]
Since $F_{i,-j}(\sigma,w)=F_{i,j+1}(\sigma,w)$ and thus
$F_{1,0}(\sigma,w)=F_{1,1}(\sigma,w)$, this may be simplified to
\begin{equation}\label{Eq_setsum}
\sum_{\varnothing\subset I\subseteq\{2,\dots,k\}}
\sum_{\sigma,w\in S_k}\sgn(\sigma)
\prod_{i=1}^k x_{w_i}^{\sigma_i}
\prod_{i\in I} F_{i,i-1}(\sigma,w)
\prod_{i\in\{1,\dots,k\}\setminus I}
\prod_{i\in I} F_{i,i}(\sigma,w)=0.
\end{equation}
Since \smash{$\prod_{i=1}^k x_{w_i}^{\sigma_i}=\prod_{i=1}^k x_i^{(\sigma w^{-1})_i}$},
we replace $w\mapsto w\sigma$ to obtain
\[
\sum_{\varnothing\subset I\subseteq\{2,\dots,k\}}
\sum_{\sigma,w\in S_k}\sgn(\sigma)
\prod_{i=1}^k x_i^{(w^{-1})_i}
\prod_{i\in I} F_{i,i-1}(\sigma,w\sigma)
\prod_{i\in\{1,\dots,k\}\setminus I}
\prod_{i\in I} F_{i,i}(\sigma,w\sigma)=0.
\]
In the following, we will show the stronger vanishing result
\begin{equation}\label{Eq_zero}
f_{I;w}:=\sum_{\sigma\in S_k}\sgn(\sigma)
\prod_{i\in I} F_{i,i-1}(\sigma,w\sigma)
\prod_{i\in\{1,\dots,k\}\setminus I}
\prod_{i\in I} F_{i,i}(\sigma,w\sigma)=0
\end{equation}
for fixed $\varnothing\subset I\subseteq\{2,\dots,k\}$ and $w\in S_k$.
For $j\in\{1,\dots,k-1\}$, let $s_j\in S_k$ denote the $j$-th adjacent
transposition, i.e., the 2-cycle $(j,j+1)$.
Then, for any $\tau\in S_k$,
\[
(\tau s_j)_i=\begin{cases}
\tau_{j+1} & \text{if $i=j$}, \\
\tau_j & \text{if $i=j+1$}, \\
\tau_i & \text{otherwise}.
\end{cases}
\]
Hence
\begin{align*}
F_{i,i-1}(\sigma s_j,w\sigma s_j)&=\begin{cases}
F_{i-1,i-1}(\sigma,w\sigma) & \text{if $i=j+1$}, \\
F_{i,i-1}(\sigma,w\sigma) & \text{if $i\neq j,j+1$}
\end{cases}
\end{align*}
and
\begin{align*}
F_{i,i}(\sigma s_j,w\sigma s_j)&=\begin{cases}
F_{i+1,i}(\sigma,w\sigma) & \text{if $i=j$}, \\
F_{i,i}(\sigma,w\sigma) & \text{if $i\neq j,j+1$}.
\end{cases}
\end{align*}
We now fix the index $j$ of $s_j$ as $j=\min\{I\}-1$.
Then $j+1\in I$ and $j\notin I$, and thus by replacing~$\sigma$
by~$\sigma s_j$ in \eqref{Eq_zero} and using that
$\sgn(\sigma s_j)=-\sgn(\sigma)$,
\begin{align*}
f_{I;w}&=-
\sum_{\sigma\in S_k}\biggl(\sgn(\sigma)
F_{j+1,j}(\sigma s_j,w\sigma s_j) F_{j,j}(\sigma s_j,w\sigma s_j) \\
& \qquad\qquad\quad\times
\prod_{\substack{i\in I\\[1pt] i\neq j+1}} F_{i,i-1}(\sigma s_j,w\sigma s_j)
\prod_{\substack{i\in\{1,\dots,k\}\setminus I \\[1pt] i\neq j}}
\prod_{i\in I} F_{i,i}(\sigma s_j,w\sigma s_j)\biggr) \\
&=-\sum_{\sigma\in S_k}\biggl(\sgn(\sigma)
F_{j,j}(\sigma,w\sigma) F_{j+1,j}(\sigma,w\sigma)
\prod_{\substack{i\in I\\[1pt] i\neq j+1}}\! F_{i,i-1}(\sigma,w\sigma)
\!\prod_{\substack{i\in\{1,\dots,k\}\setminus I \\[1pt] i\neq j}}
\prod_{i\in I} F_{i,i}(\sigma,w\sigma)\biggr) \\
&=-\sum_{\sigma\in S_k}\biggl(\sgn(\sigma)
\prod_{i\in I} F_{i,i-1}(\sigma,w\sigma)
\prod_{i\in\{1,\dots,k\}\setminus I}
\prod_{i\in I} F_{i,i}(\sigma,w\sigma)\biggr) \\
&=-f_{I;w}.
\end{align*}
Hence $f_{I;w}=0$.

Next we consider \eqref{Eq_doteddote2}.
Proceeding as in the first proof, this time it suffices to
prove the formal identity
\begin{align*}
&\sum_{w\in S_k} \det_{1\leq i,j\leq k}\bigl(x_{w_i}^j
\bigl(E_{z_{w_i}+i-j}-E_{z_{w_i}-i-j}\bigr)\bigr) \\
&\qquad{}=\sum_{w\in S_k}\det_{1\leq i,j\leq k}\bigl(
x_{w_i}^j \bigl(E_{z_{w_i}+i-j-1}+E_{z_{w_i}+i-j}-
E_{z_{w_i}-i-j}-E_{z_{w_i}-i-j+1}\bigr)\bigr).
\end{align*}
This is the same as
\begin{equation}\label{Eq_Ftilde}
\sum_{w,\sigma\in S_k}\sgn(\sigma)
\prod_{i=1}^k x_{w_i}^{\sigma_i} \tilde{F}_{i,i}(\sigma,w)
=\sum_{w,\sigma\in S_k}\sgn(\sigma) \prod_{i=1}^k
x_{w_i}^{\sigma_i} \bigl( \tilde{F}_{i,i-1}(\sigma,w)+
\tilde{F}_{i,i}(\sigma,w)\bigr),
\end{equation}
where
\[
\tilde{F}_{i,j}(\sigma,w):=E_{z_{w_i}+j-\sigma_i}-E_{z_{w_i}-j-\sigma_i}.
\]
Since $\tilde{F}_{i,0}=0$ and thus $F_{1,0}=0$, \eqref{Eq_Ftilde}
can be rewritten exactly as \eqref{Eq_setsum} but with $F$ replaced
by $\tilde{F}$.
The remainder of the proof if identical to the first proof.
\end{proof}

\end{document}
